\begin{chapterbox}
    \chapter{Tion erzählt von Mutter Natur (2021)}
    \label{Tion erzählt von Mutter Natur (2021)}
    \az{Jahr 67}
    Novize Komu vom Baum der Lieder fragt Bewahrer Tion nach den Urzeiten des Andorversums. Tion kann diesem riesigen Thema natürlich niemals gerecht werden, will es aber nicht lassen, eine Geschichte über seine liebste Göttin vorzutragen.
\end{chapterbox}

\az{Jahr 67}

\textit{"Wie ein Gewitter, das man fürchtet und gleichzeitig bewundert."}

\textit{"Sie sind die Seele des Flusses und erhalten ihn am Leben."}

\textit{"Ein magisches Wesen, das als der Urvater aller Trolle [gilt]."}

\textit{"[...]dem wohl mächtigsten Wesen der ganzen südlichen Welt."}

\textit{"Die Trolle lebten schon immer in Andor."}

\textit{"Die Zwerge [gruben] tief unter der Erde Schächte und riesige Hallen, die den Drachen als Heimstatt dienten. Im Gegenzug übergaben die Drachen den Zwergen die Geheimnisse des Feuers und erlaubten ihnen damit, die ewige Düsternis unter dem Gebirge zu vertreiben."}

\textit{"Krahal [...] wurde vor vielen Hundert Jahren tief erschüttert."}

\textit{"Was ich sah, war ein Geflecht aus Düsternis, Umrisse eines Baumes, in dessen Adern rotes Blut floss."}

\textit{"Die letzten Drachen allerdings starben oder versteinerten mit den Jahren. Nur bei einem von ihnen, dem jungen Drachen, der die Kraft des Feuerschildes mit eigenen Augen gesehen hatte, erlosch die Flamme des Zorns nie."}

\textit{"Sie beteten leise zu Mutter Natur, sie möge dem Kranken Kraft schenken und ihn gesunden lassen."}

\textit{"Dort, glaubten die Andori, würde das ewige Glück auf sie warten."}\bigskip

Novize Komu vom Baum der Lieder leerte derart viele Schriftrollen und Schriftstück-Fetzen auf einmal aus der Holzkiste, dass einige davon fröhlich vom Schreibpult rollten und auf den Stammboden kullerten. Während Komu hastig auf den Boden sank und die Rollen von Weiterrollen abhielt, legte Meister Tion seinen Kopf schief und zog seine buschigen Augenbrauen zu einem fragenden Blick zusammen.

"Warum so hektisch, Komu? Liegt dir etwas auf dem Herzen?"

Komu stockte einige Male, aber sobald der Redeschwall einmal begonnen hatte, war dieser kaum zu bremsen. Enthusiasmus für die Arbeit der Bewahrer hatte das Kind, das musste man ihm lassen.

"Nun, ich habe mich kürzlich gefragt... wir haben ja in der Schule viele... viele Bruchstücke aus der reichen Vergangenheit Andors haben wir schon vernommen, aber... nun, es ist so... die Stücke der Geschichten sind da, aber zwischen ihnen liegen immer noch große Lücken, und ich habe Ega darauf angesprochen. Ega meinte, dass wir noch ganze zwei Jahre warten müssten, bis wir soweit wären, einen ersten geordneten Überblick über die Urzeiten zu erhalten... dass das anspruchsvolle Texte sind, das verstehe ich ja, aber so lange warten mag ich nicht, und da... da fragte ich mich... nun... wie passen diese Stücke alle zusammen? Wie lautet denn die ganze Geschichte der Urzeit? So, nur in der Kurzfassung. Also, falls man das schon verstehen kann mit unseren Kenntnissen."

Tion legte seine Schreibfeder beiseite und kratzte seinen langen Bart. Dann schmunzelte er.

"Ach, Komu, eine ‘ganze’ Geschichte gibt es nicht, das wirst du wahrscheinlich bald verstehen. Wir bewahren bloß Schriften aus den vergangenen Zeiten und fertigen neue Aufzeichnungen an. Je länger deine Ausbildung sich hinziehen wird, desto mehr Schriftstücke wirst du kennen lernen und miteinander in Verbindung setzen können. Dabei wirst du aber auch feststellen, dass es nicht den ‘einen’ Verlauf einer Legende gibt, sondern ganz unterschiedliche Berichte, welche zum Teil gar widersprüchlich sind. Es ist eine Kunst, einen möglichst objektiven Blick auf das Vergangene anzustreben, eine Kunst, die man wie alle Künste nur durch langes Üben meistern kann."

Tion wandte sich schon wieder seiner Schreibfeder zu, da hielt er inne und murmelte: "Was soll‘s, ich könnte ohnehin etwas Ablenkung gebrauchen. Hast du eine spezifische Frage oder Geschichte im Sinn, zu der ich ausführen könnte?"

Komu überlegte kurz und plauderte dann fröhlich los: "Hat Mutter Natur uns nun erschaffen oder nicht? Weil, meine Eltern sagten, dass das so sei, aber Ega behauptete, dass in den alten Texten was ganz anderes steht... und wie kam es eigentlich zum Streit zwischen Mutter Natur und dem Urtroll? Weil, so, wie ich das gehört habe, war der Urtroll am Anfang doch ein ganz gutmütiger Geselle, und jetzt wird er als so gefährlicher Gegner beschrieben, und irgendwie hängt doch alles damit zusammen, dass der Baum der Lieder überhaupt..."

"Irgendwie hängt doch alles mit allem zusammen", unterbrach Tion Komus Redeschwall zwinkernd, dann fuhr er fort, "Mutter Natur behüte deinen Eifer, junger Novize. Ich werde dir natürlich keine eindeutigen Antworten geben können. Wie viel Wirken Mutter Natur in unserer Entstehung hatte, ist debattierbar, auch wenn ihr Einfluss insbesondere auf unseren Orden sicherlich gigantisch war. Und ich könnte vielleicht eines Tages die Geschichte mit dem Urtroll erzählen, so wie ich sie in Erinnerung habe. Aber wie viel Wahrheit noch dahinter steckt, musst du am Ende selbst beurteilen. Es ist schon eine Zeit lang her, dass ich die Urtexte studiert habe. Das Niederschreiben der aktuellen Ereignisse im Lande nimmt einige Aufmerksamkeit in Anspruch."

Komu war fröhlich aufgehüpft, sobald Tion eine mögliche Geschichtenerzählung angesprochen hatte. Tion konnte nicht anders als zu grinsen und brummelte: "Nun gut, jetzt hast du mich mit deinem Fieber angesteckt. Trommle uns ein paar junge Adepten zusammen, ich bin in Erzähllaune."

Komu stieß eine Faust in die Höhe und trabte davon, um geschwind wie der Wind Freunde zu rufen. Tion blieb zurück und rief sich die wichtigsten Punkte zur Geschichte von Mutter Natur und dem Urtroll in Erinnerung. Auch wenn er es nicht direkt zugeben würde, mochte er es ungemein, dem Nachwuchs die alten Geschichten näher zu bringen. Würde er es dieses eine Mal schaffen, nicht allzu stark abzuschweifen? Wahrscheinlich nicht. Aber das war auch nicht die Hauptsache. Hmmm... wo wäre ein guter Einstieg für seine Erzählung?\bigskip

"Vom Helden Fenn haben wir erfahren, was die wilden Barbaren des Ostens ihren Kindern erzählen. Dass diese Welt einst leer und kahl gewesen sei bis auf die drei Götter, welche sie Stück für Stück mit Leben füllten – manche davon kreative Neukreationen wie die Menschen und Büffel, andere bloß Mischungen ihrer vorherigen Kreationen, etwa die Nixen. Die Götter ließen ihre Geschöpfe gegeneinander antreten, doch dann zogen sie sich immer mehr zurück, während ihr Interesse nachließ, bis sie bloß noch mit den wenigen Schamanen der Barbaren-Stämme kommunizierten, und auch das immer seltener.\bigskip

Manche Schildzwerge glauben, dass sie einst von einer Urmacht geschaffen wurden, um in deren Namen der Erde ihre Geheimnisse zu entlocken. Sie erzählen Legenden von einer utopischer Unterwelt tief unter der unseren, voller Gold, glitzernder Edelsteine und magischer Schätze, aber auch voller Feuer, Gefahren und finsteren Kreaturen. Sie graben selbst heute noch im Auftrag dieser Urgewalt tiefer und tiefer in die Tiefminen. Viele hoffen auf einen Fingerzeig dieser Entität. Seit einigen Jahrzehnten wird hier und da die Behauptung aufgerufen, dass der Tod Jari Dorrs und die Entdeckung der Silbermine ein solches Zeichen gewesen sei, das die Zwerge in den Norden riefe, doch die meisten älteren Schildzwerge tun dies als Humbug ab.\bigskip

Eine Menge Zauberer aus der fernen Eiswelt Hadrias meinen, dass höchstwahrscheinlich niemand eine Hand in ihrer Entstehung gehabt habe – höchstens eine Faustvoll wilder Magie, die aus der Hadrischen Unterwelt wie Dampf aufsteigt und alles Leben zumindest ein wenig formt.\bigskip

Nun, wir, die Bewahrer vom Baum der Lieder, glauben an das, was unsere Vorfahren lange vor unserer Zeit erlebten und niederschrieben. Diese Schriftstücke wurden in diesem Baum über die Jahrhunderte hinweg bewahrt, und viele detaillieren die Geschichte von Mutter Natur.

Die ersten Bewahrer beschreiben Mutter Natur als das erste Wesen, das das Licht der Welt erblickte und durch diese Landschaften strich. Manche waren sich sogar sicher, dass sie es gewesen sei, die dieses Licht der Welt überhaupt erst erschaffen habe. Fest steht, dass Mutter Natur sich seit Anbeginn der Flüsse, Berge und Wälder um die Natur sorgte und sich daran erfreute, wie die Sonne das Land beschien, die nassen Felsen in ihrem Licht glitzerten und allerlei Getier durch das Gras wuselte. Und da Freude durchs Teilen nur verstärkt wird, beseelte Mutter Natur das Land mit ihren Kindern. So schlüpften die Wassergeister in die Flüsse, die Erdgeister ins Fahle Gebirge, die Waldgeister in die Bäume und die Geister des Windes und der Wolken, angeführt vom edlen Arkteron, in die Lüfte. Die Naturgeister waren weitaus gescheiter als die Tiere des Landes, da Mutter Natur in jeden von ihnen einen Kern ihrer selbst pflanzte. Und Mutter Natur erfreute sich an den Seelen des Landes und hegte und pflegte sie so wie die restliche Natur.

Doch als Mutter Natur die Naturgeister erschaffen hatte, so konnte – oder wollte – sie sie nicht länger kontrollieren. Und während viele Geister damit glücklich waren, einfach zu existieren, zu interagieren und gemeinsam innerhalb der Ordnung dieser Welt umherzuspielen, so kam es, dass manche Geister sich von Mutter Naturs Führung abwandten und sich der Zerstörung zuschrieben. Diese empfanden nur Freude daran, anderen die Freude zu nehmen. Fürchterliche Feuergeister setzen die Bäume in den Brand, verdampften das Wasser, vertrockneten die Erde und verbrannten die Pflanzen.

Da war Mutter Natur traurig und versuchte, mit ihnen zu reden. Einige Feuergeister baten um Entschuldigung, und die ward ihnen gewährt. So entstanden die Feuertakuri des Westens.

Doch andere abtrünnige Geister, des Feuers, des Schattens und der dunklen Erde, wollten ihr Verhalten nicht ändern. Sie fürchteten, dass Mutter Natur sie richten würde, und zogen davon, verließen Mutters Lande und suchten sich ihr eigenes Heim, weit im Süden, wo Mutters Einfluss immer schwächer wurde. Die Erdgeister unter ihnen falteten das Land und schichteten das Graue Gebirge auf zwischen dem Südland, ihrer neuen Heimat, und dem Norderland von Mutter Natur. Sie hatten die Hoffnung, dass das Graue Gebirge sie vor Mutter Natur schützen würde, falls diese je beschließen sollte, etwas gegen ihre Zerstörung zu unternehmen.

Doch in ihrem eigenen Reich waren die abtrünnigen Geister immer noch von Sehnsucht nach ihrer alten Heimat erfüllt. Darum versuchten sie, diese zu reproduzieren. Sie erschufen Flüsse, doch nichts als Lava floss in ihnen. Sie formten Bäume, doch zerfielen diese zu Staub und Dreck. Sie erschufen ihre eigenen Tiere, doch diese klappten kraftlos und kalt zusammen, denn die Naturgeister kannten das Geheimnis des Lebens nicht.

Auch wenn keine schönen Kreationen dabei zustande kamen, so zogen die Taten der Geister die Aufmerksamkeit eines mächtigeren Wesens auf sich, das durch die Südlande streifte. Es war der Urtroll, und er kannte das Geheimnis des Lebens. Bislang hatte er sich nicht groß um diese Welt geschert. Doch als er sah, was die abtrünnigen Naturgeister zu schaffen versuchten, da fühlte der Urtroll Mitleid mit ihnen, und er hauchte ihren Kreationen Leben ein. Die entstandenen Wichte waren krank und missgestaltet, und alsbald krochen sie in die Höhlen unter dem Lande und versteckten sich vor dem Anblick derjenigen, die direkt Mutters Güte entstammten.

Da wurde der Urtroll dazu inspiriert, seine eigenen Wesen zu schaffen. Er formte aus der Erde des Grauen Gebirges die Trolle, und hauchte ihnen Leben ein, und die Trolle breiteten sich in die umliegenden Lande aus und richteten Verwüstung an. Mutter Natur hatte schon einige Zeit lang mit Sorge die Handlungen der abtrünnigen Geister betrachtet, und nun hielt sie es für nötig, einzugreifen, um ihre Welt vor den Trollen zu beschützen. Doch sah sie auch, dass der Urtroll stolz war auf seine Kinder, und dass man mit ihm reden konnte. So beschloss Mutter Natur, mit dem Urtroll verhandeln. Sie wollte allerdings nicht selbst in die Südlande reisen und Furcht und Schrecken unter den abtrünnigen Naturgeistern verbreiten, denn sie hatte auch ihre entflohenen Kinder immer noch lieb. So suchte sie stattdessen nach Boten, die ihre Nachricht zum Urtroll bringen könnten.

In der Zwischenzeit hatte Mutter Natur weiterhin die Tiere, Pflanzen und Waldpilze in ihrem Reich gehegt und gepflegt, und fröhlich erkannte sie, dass viele davon durch ihre gute Pflege selbst Seelen entwickelt hatten, ja einige gar in rudimentärem Kontakt zu den friedlichen Naturgeistern des Landes standen. Unter diesen vielen Arten fiel Mutter Natur besonders eine ins Auge, eine Sorte von breitwüchsigen Wesen, welche tief unter der Erde nach Geheimnissen suchten. Mutter Natur gab ihnen den Namen Zwerge und schenkte ihnen die Magie der Runen. Sie erschien dem Zwerg Eibert und sang ihm eine sanfte, klare Melodie vor, die ihn mit Kraft und Mut erfüllte. Mutter Natur bat Eibert, diese Töne mit seiner Zwergenschar zu teilen, sich durch das Graue Gebirge zu graben und das Lied dem Urtroll und dessen Brut vorzuspielen. Dies würde ihn hoffentlich besänftigen, sodass er der Zerstörung der Welt durch seine Kinder gewahr werden und eingreifen möge.

Eibert nannte die Melodie Steinsang, und er tat, wie von Mutter Natur gebeten. Doch kaum hatten die Zwerge mit den Grabungen begonnen und erste Konflikte mit den Erdgeistern ausgetragen, da stürzten die Drachen aus dem Himmel herab. Das waren mächtige magische Wesen, die vom roten Mond stammten und die Kunst des Feuers beherrschten. Manche unter ihnen waren wie die abtrünnigen Naturgeister dem Durst der Zerstörung erlegen und jagten bösartig durch das Land, doch andere von ihnen waren gutherzig und suchten die Nähe der restlichen Wesen.

So kam es, dass die Zwerge von Mutter Naturs Auftrag abkamen, um sich stattdessen um die Plage der bösartigen Drachen zu kümmern. Sie schlossen Frieden mit den gutmütigen Drachen und bauten ihnen große Hallen unter der Erde. Im Gegenzug teilten die Drachen mit ihnen das Geheimnis des Feuers. Die Zwerge nutzen das Feuer, um tiefer denn je unter die Erde zu dringen, tiefer als alle anderen Gänge Caverns, und dort schlugen sie eine gigantische Höhle in den Fels, groß genug für hunderte von Drachen. Die gutmütigen Drachen versammelten sich und ließen ihre ganze Magie in diese Höhle fließen, und da wuchs aus ihr ein riesiger, weiß schimmernder Baum, durch dessen kräftige Adern stetig rotes Blut gepumpt wurde. Ein lebendiges Zeichen des Friedens.

Dieser magische Ort war Krahal. Die Geister der Drachen vereinten sich in Krahal und wurden eins. Die Güte der Drachen, die den Ort gegründet hatten, umschloss den Zorn der bösartigen Drachen und verschloss ihn im weißen Baum. Das Böse war zumindest für eine Zeit lang nicht mehr Teil der Drachen. Ein Frieden zwischen Drachen und Zwergen wurde geschlossen. Anstatt sich auf Mutter Naturs friedlichen Auftrag zu besinnen und Steinsang zum Urtroll zu bringen, spielten die Zwerge und Drachen Steinsang nur für sich selbst, um sich mit noch mehr Mut und Kraft zu erfüllen. So drohten sie den Trollen.

Mutter Natur hatte mit Freude beobachtet, wie die Drachen sich auf die Seite ihrer Welt stellten, doch sah sie nun mit Sorge, wie der Urtroll um seine Kinder fürchtete und sich noch mehr von ihr zurückzog. So suchte sie nach einem neuen Botschafter. Ein Schattengeist trat ins Licht und bot an, Mutter Naturs Bitten zum Urtroll zu tragen. Mutter Natur konnte ihm nicht vertrauen, doch wie sich der Konflikt zwischen den Zwergen, Drachen und Trollen stetig zuspitzte, drängte die Zeit. Darum übergab sie dem Schattengeist eine neue Melodie und bat ihn, diese zum Urtroll zu bringen und ihm vorzuspielen. Diese Melodie sollte ihm ihren Blick auf die Welt zeigen und ihn einen Frieden zwischen den Völkern wertschätzen lassen.

Doch Mutter Naturs Botschafter war ein gerissener Geist. Es war ein Schattengeist, der sich innerlich schon längst von Mutter Natur abgewandt hatte, doch nicht mit den restlichen abtrünnigen Naturgeistern in den Süden geflohen war. Es war der Schwarze Herold, und nun erkannte er seine Gelegenheit, großes Leid anzurichten, und freudig tat er es. Er verdrehte die Melodie aus dem Mund von Mutter Natur und als er die dadurch korrumpierten Töne dem Urtroll vorsang, so wiegelte er den Urtroll gegen Mutter Natur auf.

Der Urtroll erhob sich brüllend vor Wut und sog die Kraft aller abtrünnigen Geister in sich ein. Er wuchs und wuchs, bis er so gigantisch war wie das Land selbst, sein Körper weit aus den Wolken ragte und sein linkes Horn gegen den roten Mond splitterte, von dem die Drachen einst gefallen waren. Der Urtroll stampfte mit seinem steinernen Fuß auf und das Graue Gebirge splitterte entdrei, mit tiefen Schluchten, die bis ins Reich der Zwerge vordrangen. Mit einem einzigen Schritt stapfte der Urtroll tief in das Herz von Mutter Naturs Reich, packte sie mit seiner gigantischen Faust und schleuderte sie auf die Erde, so heftig, dass ihr Land sich unter der Schockwelle in tausend Stücke teilte. Diese tausend Stücke drifteten vom Grauen Gebirge davon und wurden zu den tausend Inseln des Hadrischen Meeres.

Mutter Natur lag auf dem letzten Stück Land, welches noch in Verbindung zum Grauen Gebirge stand, und lag im Sterben. Doch selbst in diesem Moment schlug sie nicht gegen den Urtroll zurück, sondern hielt ihm ihre uneingeschränkte Liebe entgegen. Zum Glück erkannte der Urtroll da, dass die misstönige Melodie, die ihn so sehr in Rage versetzt hatte, nie von so einem puren Wesen hätten stammen können, und er verfluchte den hinterhältigen Schwarzen Herold, der ihn in Rage versetzt hatte.

Voller Reue versuchte der Urtroll, Mutter Natur zu retten. Er ließ seine Macht in sie fließen, und er schrumpfte und schrumpfte, bis er kaum grösser war als die größten Trolle, die er damals geformt hatte. So rettete der Urtroll ihr Leben, doch seine Tat des Hasses konnte nicht so einfach ungeschehen gemacht werden. Mutter Natur schlug ihre Augen nicht mehr auf. So sank der Urtroll neben Mutter Natur auf den Boden, bettete sie auf ein Bett aus Kräutern und sang ihr mit seiner urtiefen Stimme eine eigene Melodie, voller Trauer und Scham, aber auch voller Kraft und Hoffnung. Und die Zwerge kamen zusammen, und spielten Steinsang auf Flöten. Die friedlichen Naturgeister traten hinzu und ließen ihre Klagelieder ertönen. Auch unsere Vorfahren kamen daher und trugen ihre eigene Musik vor. Und alle diese Lieder vereinten sich zu einem riesigen Gebilde. Entgegen aller Erwartung überlagerten sie sich zu einem wunderschönen Klang und trieben jedem Zuhörer die Tränen in die Augen. Das war das Schlaflied von Mutter Natur.

Das Land summte und brummte mit der Kraft, die das Schlaflied von Mutter Natur mit sich trug. Dort, wo Mutter Natur auf einem Bett aus Kräutern schlief, wuchs ein prächtiger Baum, grösser und kräftiger als alle anderen des Waldes, und vollkommen von der Magie der Lieder erfüllt. Erst als seine Wurzeln Mutter Natur vollkommen umschlossen hatten, verklang die Musik der vereinigten Völker und sie gedachten ihr.

Unsere Vorfahren blieben zurück und versuchten, das Gehörte zu Pergament zu bringen. Es sollte ihnen nicht gelingen, den perfekten Klang des Schlafliedes auch nur im Ansatz zu reproduzieren, doch hörten sie nie mit dem Niederschreiben von Liedern auf. Als ihnen die Lieder auszugehen begannen, so notierten sie sich die Geschichte des Landes, und als auch diese immer ausgeschöpfter wurde, wandten sie sich schlussendlich den Legenden zu. Sie besiedelten den mächtigen Baum über Mutter Naturs Schlafstätte und füllten ihn mit Liedern und Wissen. Mit der Zeit sollten aus ihnen die Bewahrer des Baumes der Lieder werden. Doch noch war es nicht an der Zeit.

Nun, da Mutter Natur im Schlafen lag, gab es natürlich niemanden mehr, der die Natur unter Kontrolle hielt. Erdbeben tobten über das Land. Stürme entwurzelten die Bäume. Blitze spalteten Stein. Es war, als ob die Erde selbst für viele Jahre rebellierte.

So kam es in der Zeit der Not, dass viele Völker sich voneinander teilten und sich nur noch um sich selbst kümmerten. Die Trolle richteten weiterhin Verwüstung an, und quälten Mensch und Tier. Die Menschen wurden zu vielen kleinen Clans, die sich in alle Himmelsrichtungen verstreuten und ihre Herkunft vergaßen. Die Zwerge zogen sich mit den Drachen ins Gebirge zurück. Manche wandten sich gar von den Drachen ab und wurden zum zurückgezogenen Gebirgsvolk der Agren. Viele Naturgeister vergingen in Trauer und wurden wild, als Mutter Naturs Präsenz sie verließ. Arkteron, der Herr der Winde, hielt es in der Nähe von Mutters Schlafstätte nicht mehr aus und zog, von Trauer und Hass erfüllt, gen Norden. Dort ward er zum Herrn der Stürme und traf auf zwei weitere mächtige Wesen. Gemeinsam ersuchten sie, das Wunder des Lebens dem Urtroll nachzumachen und ihre eigenen Trolle zu erschaffen. Doch das ist eine Geschichte für ein anderes Mal.

Tiefe Schluchten zogen sich seit dem Anfall des Urtrolls durch das Graue Gebirge. Eine von ihnen wurde im Laufe der Zeit zu einem gigantischen Drachenhort, in dem sie ihre Körper ruhen ließen. Krahal-Schlucht wurde sie genannt. Eine andere hingegen wurde zu einem steinernen Mahnmal der Geschichte, denn der Urtroll selbst legte sich in dieser Schlucht nieder, als er von Scham und Zorn erfüllt ins Graue Gebirge zurückkehrte. Er wollte nur noch schlafen und nicht mehr wissen, zu welchem Unheil der Schwarze Herold ihn angestiftet hatte, und so bat er die Urahnin aller Agren, die alte Korn, ihn alles vergessen zu lassen.

Korn fühlte mit dem Urtroll und wirkte gemeinsam mit den Druiden der Agren einen Zauber, der den Urtroll in einen tiefen steinernen Schlaf sinken ließ. Sie ließ aber eine kleine Knochenflöte anfertigen, die den Urtroll wieder ins Reich der Wachen zurückholen würde, sollten die Agren je auf seine Hilfe angewiesen sein. Fortan wachten Korn und ihre Nachfahren über die Knochenflöte und die Schlucht, welche bald als Korn-Schlucht bekannt wurde.

Nur die abtrünnigen Naturgeister des Südens freuten sich ob des Machtvakuums im Lande. Sowohl Mutter Natur als auch der Urtroll schliefen nun, und so trat der Schwarze Herold erneut hervor. Er hatte beobachtet, wie der Urtroll die Trolle erschaffen hatte. Er kannte jetzt das Geheimnis des Lebens. Und er hatte vor, ein Volk zu erschaffen, welches in seinem Namen die Welt knechten solle. Doch alleine hatte er nicht die Kraft dazu. So versammelte der Schwarze Herold weitere Schattengeister, Erdgeister und Feuergeister, und er schöpfte aus ihrer Macht. Gemeinsam hatten sie die Mittel, eine Spezies zu erschaffen, größer und stärker als alle anderen Humanoiden, ohne Skrupel und ohne Gewissen. Das waren die Riesen, die das öde Reich des Südens für sich beanspruchten.

Doch konnte bald jeder sehen, dass die Riesen den Drachen bei weitem unterlegen waren, und dass sie dem Schwarzen Herold nicht gehorchten. Sie wandten sich von ihm ab und ihren eigenen Herrschern zu. Dunkle Hexerei entdeckten manche, und die Krahder nannten sie sich fortan. Nach und nach verloren immer mehr abtrünnige Naturgeister das Vertrauen in die Führung des Herolds, und da verließen sie ihn und zogen alleine durch eine Welt, welche immer geteilter und verrückter wurde. Selbst die Drachen und die Zwerge verwickelten sich in einen Unterirdischen Krieg, der das Graue Gebirge erschütterte und beide Völker an den Rande der Vernichtung bringen sollte.

Der Schwarze Herold blieb alleine zurück und hatte einen letzten Plan, die Anerkennung seiner Kreationen wieder zu erringen. Es zog ihn nach Krahal, den magischen Ort, den die Drachen erschaffen hatten. Krahal war erfüllt von Wut und Verzweiflung ob des wütenden Krieges, und der darin stehende weiße Baum verfaulte langsam unter den Qualen des Drachenvolks. Der Schwarze Herold erhoffte noch immer, das Reich des Südens zu einem Abbild von Mutter Naturs Reich zu machen, und so ergriff er eine Schwarze Wurzel des verfaulten Baumes, brach sie ab und brachte sie zu Nomion, einem gerissenen Krahder-Hexer. Dieser solle sie nähren, aus ihren einen mächtigen Schwarzen Baum wachsen lassen, daraus Kraft ziehen und dafür den Schwarzen Herold anbeten.

Doch wie schon der Schwarze Herold zu Mutter Natur wandte sich auch hier die Kreation von ihrem Kreator ab. Nomion ließ seine Schüler die Schwarze Wurzel aus Krahal in Krahd einpflanzen, und daraus wuchs der Schwarze Baum der Dunklen Hexerei. Nomion nutzte die Kraft des Baumes, um das ewige Leben des Schwarzen Herolds zu rauben. Der Schwarze Herold wurde zu einem Schatten seiner selbst degradiert, nur noch eine leere Hülle, welche wie ein Unwetter durch die Lande strich und das Böse unterstützte, wo immer sie konnte. Der Schwarze Herold verfluchte Mutter Natur dafür, ihn erschaffen zu haben, und richtete seine Wut auf sie.\bigskip

Hier enden die Berichte und Legenden unserer frühsten Vorfahren. Doch können wir natürlich folgern, was mit dem Schwarzen Herold geschah, während die Zeit verging.

Während der Unterirdische Krieg zwischen Zwergen und Drachen tobte, stahl Nomion Korns Knochenflöte und weckte den Urtroll. Dieser hatte in seinem tiefen, traumlosen Schlaf alles vergessen bis auf seine Wut, und so war er kaum mehr eine Urgewalt, die ohne Sinn und Verstand tobte und zerstörte. Eine Urgewalt, die Nomion mit Dunkler Hexerei zu lenken vermochte. Wäre Nomions Körper nicht kurz darauf vernichtet worden, wer wüsste, was er alles hätte anrichten können. Nur durch das ewige Leben des Herolds, das Nomion diesem durch den Schwarzen Baum gestohlen hatte, konnte ein Teil seiner Seele fortbestehen. Aber das ist eine Geschichte für ein anderes Mal.

Auch ohne Nomion und den Urtroll zogen die Krahder aus und raubten Menschen, Zwerge und Trolle aus den umgebenden Ländern, um sie sich selbst untertan zu machen. Der Schwarze Herold setzte seine Hoffnung in die Könige der Krahder, doch all diese wagten nicht, weiter gen Norden zu ziehen und Mutter Naturs Baum der Lieder zu verbrennen, denn auf dem Weg wachten noch immer die letzten Drachen. Nicht mehr viele von ihnen waren seit dem Verklingen des Unterirdischen Krieges noch übrig, und mit der Zeit waren auch von den Überlebenden viele gestorben oder versteinert. Doch noch immer waren sie eine zu große Gefahr, als dass die Krahder es wagen würden, sich ihnen direkt zu stellen.

Unter den letzten Drachen war Tarok, der stärkste der Drachen, der den Ausbruch des Krieges miterlebt hatte und vor Zorn schon beinahe wahnsinnig geworden war. So setzte der Schwarze Herold Hoffnung in Tarok und zeigte ihm, wie der Urtroll einst die Trolle erschaffen hatte. Tarok folgte den Weisungen seines Herolds. So formte Tarok in Krahal die Körper der gefallenen Drachen zu wüsten Geschöpfen, unsagbare Kreaturen mit zerstörerischen Kräften, in Anlehnung an die Zwerge und Menschen der Umgebung. Sie sollten in seinem Auftrag ins Zwergenreich einfallen und den Ausgang des Unterirdischen Krieg doch noch an die Drachen reißen.

Dazu kam es nicht. Taroks unsägliche Taten erschütterten den magischen Ort zutiefst. Rotes Blut spritzte aus dem verfaulten Baum Krahals. Die Geister der restlichen Drachen vergingen und ihre Körper stürzten seelenlos zu Boden. Einzig der wahnsinnige Tarok, für den Angst und Schmerz, Trauer und Verzweiflung nunmehin zum Lebenselixier geworden waren, konnte sich mit schierer Willenskraft an diese Welt klammern und die Erschütterung Krahals überstehen. Fortan musste er damit klarkommen, dass er es gewesen war, der den Untergang der letzten Drachen ausgelöst hatte, und der letzte Funke Freude in seinem Herzen erlag dem Schmerz.

Taroks Kreaturen verließen Krahal und scharten sich in den Bergen und Wäldern von Andor, um für den letzten Drachen zu kämpfen und zu töten. Und wie sie Krahal verließen, so drang Taroks Geist in die Trolle ein und auch sie folgten fortan seinem bösartigen Befehl. Und der schwarze Herold diente Tarok und trieb seine Kreaturen an.

Doch wie wir alle wissen, sollte Tarok sich nicht an die Wünsche des Schwarzen Herolds halten und nicht den Baum der Lieder fällen. Der letzte Drache wurde von seinem Rachedurst gegen König Brandur von Andor abgelenkt und ließ von den Plänen des Herolds ab. Am Ende wurde er von den Helden von Andor überwunden. Die Überreste seines Körper wurden die Narne hinuntergespült, und auch wenn seine Kreaturen noch immer hasserfüllt unser Land belagern, so hoffen wir doch, dass selbst Tarok im Tode das offene Meer erreichen konnte und, nachdem er dort seine bösartigen Taten besehen und berichtig hat, wie lange es dafür auch brauchen möge, im ewigen Glück seinen Frieden finden kann.

Der Wunsch des Schwarzen Herolds ging nicht in Erfüllung. Das Böse in Andor mehrt sich nicht, sondern wird stetig schwächer. Stetig mehrt sich hingegen Mutter Naturs Macht, selbst im Traum vermag sie die Geschicke dieser Welt schon wieder zu beeinflussen und uns Kraft zu verleihen. Und wie Mutter Natur sich in ihrem Schlafe heilt, so heilt auch das Land von dem Übel der abtrünnigen Naturgeister. Eines Tages wird auch der Schwarze Herold seine Verfehlungen einsehen und auf die Seite des Lichts treten. Und eines Tages wird sich so viel Frieden verbreitet haben, dass Mutter Natur vollends erwachen und ihre Rolle als Hüterin des Landes wieder aufnehmen kann.

Wir wissen nicht, ob sie in einem Jahrzehnt erwachen wird, oder in einem Jahrhundert, oder gar in einem Jahrtausend. Aber wir glauben fest daran, dass Mitglieder unseres Ordens sie erwarten werden, wenn sie sich wieder erhebt. Bis dahin verweilen wir Bewahrer vom Baum der Lieder stets hier im Wachsamen Wald und schützen ihre Ruhestatt. Wir wahren alles Wissen, das wir können, ohne groß einzugreifen. Wie sie es uns gelehrt hatte. Und wir können es kaum erwarten, ihr berichten zu dürfen, was alles während ihres Schlafes vorfiel."
